\section{Discussion}
\label{sec:discussion}

\subsection{What the evidence supports}
The second confirmation supports a specific and bounded claim. On 35 rivers that no earlier round had used, measured at byte level, spending bytes on witnesses of hard cases decided 0.18 to 0.23 more hard cases at unanticipated thresholds than a node that sends the configured answers and also uses its record (H2), and 0.20 to 0.28 more than the record alone (H3), at a cost of 0.03 to 0.08 in the share of days whose whole record was recovered (H4). No certain decision was wrong (H1). The study does not claim that witnesses are the best way to decide thresholds that are fixed when the node is configured: the flags were better there in development and in both confirmations (H6).

The practical reading is a choice that depends on how fixed the thresholds are:
\begin{itemize}
  \item If the thresholds are known and stable, a node that sends the configured answers with its record decides nearly every configured hard case (0.999 under independent loss at six messages per day) at no cost in record recovery, because on this frame the four flag bits fit in spare bits.
  \item If thresholds may change, or if the agency needs to decide thresholds it did not configure, the targeted witness design decides the most hard cases at unanticipated thresholds (0.798 against 0.572 for flags plus record and 0.523 for the record alone), at a cost of about 0.08 in whole-record recovery at six messages per day and 0.03 at twelve.
\end{itemize}
This is a claim about use, not a measured benefit in the field: the unanticipated threshold sets were chosen by the analyst, and whether thresholds change in practice was not studied.

\subsection{What changed between the two confirmations, and why it matters}
The first confirmation failed its integrity hypothesis and, read critically, overstated the benefit of witnesses: its flags comparator discarded the record; therefore, the contrast against flags measured the record as much as the witnesses. The revision replaced the comparator by one that uses the record, moved every hypothesis to the byte level, and lengthened the check. The benefit over the fair comparator (H2) is smaller than the earlier contrast against flags (0.53 at six messages per day under independent loss) but it is the contrast that isolates the witnesses, and it held on two development years and on the sealed cohort.

Two decisions taken between the confirmations weaken the evidence and are stated here rather than in a footnote:
\begin{itemize}
  \item The 3b rule was adopted although it missed its own record clause: the adoption rule, written before the comparison, allowed a record cost of at most 0.08, and the cost on the CRC-24 frame was about 0.09 at six messages per day. The design was adopted because no other design met the benefit clause at an acceptable cost, and tested the cost (H4) with a looser limit instead. H4 held, but a reader who applies the original rule would reject the design at six messages per day under independent loss (point cost 0.082 in the cohort). At twelve messages per day the cost was 0.03 and met the rule.
  \item The first argument for that adoption, that the cost was small against the fair comparator, rested on a modelling error that an internal reading found before the freeze: the flags had been charged a record value per message that the frame does not need. The argument was withdrawn, and the hypothesis built on it was dropped. The error did not affect the benefit hypotheses.
\end{itemize}

\subsection{Integrity}
With a 24-bit check on every message, the expected rate of days with a false certificate at an assumed bit-error rate of $10^{-3}$ was at most 0.0035 per million days, and no receiver could exceed 0.084. This bound, not a test, replaces the conditional hypothesis that failed in the first confirmation. The conditional picture did not change: given a forged message that passes the check, 0.11 to 0.22 of the days on which no surplus message arrived received a false certificate, and about 2\% of forged days at six messages per day received a wrong decision. The safety therefore comes from the length of the check under a random-corruption model, not from the receiver rules alone; a structured fault or an adversary would need authentication, which is out of scope.

\subsection{Limits}
\begin{itemize}
  \item Simulation only: the evidence is trace-driven replay of USGS records through modelled loss. No satellite hardware, no measured channel error statistics, and no microcontroller measurements were used.
  \item Scope: 35 rivers of the United States in one year (2022), with one or two sites in three of the nine regions; nothing is claimed for other regions, years, or deployments.
  \item What was tested: three hypotheses could fail. H1 cannot fail through assumption R, H8 compares two implementations of the same code, H6 holds largely by construction, and H7 cannot exceed its ceiling.
  \item Assumption R: the reporting resolution was inferred from the analysed record for 26 of the 35 sites. H1 is therefore a test of the code, conditional on R.
  \item Comparator: the flags-plus-record design takes its configured answers from the true record whenever a message arrives, which makes it an idealised node; its flag field is assumed to occupy four of the 16 spare bits, and the frame does not implement it.
  \item Not re-run: send-on-delta, the tightened tolerance, the control for targeting, the exact allocator, and the hybrids were compared only in development and in the first confirmation; the final method was not tested against them.
  \item Censored codes: a forged record value with a censored code passes the sanity rule and removes a reading from decisions.
  \item Analyst choices: the alternative threshold sets, the hard-case bands, the 7DADM band, and the unknown-history rule were chosen by the analyst and were seen on every development cohort.
  \item Truth: truth statements refer to the loader's record, which includes readings bridged over gaps of at most two readings, and not to the river.
  \item Witness decisions are a lower bound on what witnesses could decide, because the box rule is not exact with witnesses (Proposition 2).
  \item Run: the confirmatory run was interrupted once by a restart of the computer and resumed with the same frozen code (Section~\ref{sec:procedure}).
  \item Review: the code, the proofs, and the plans were read by separate AI sessions, which is an in-project reading and not an outside review, and by the project team; no outside reviewer has read them.
  \item Priority: the novelty statement of Section~\ref{sec:related} rests on a general web search and on sources read in part; it is not a systematic review.
\end{itemize}
